% Working draft for MNRAS.
% Canonical copy: https://www.overleaf.com/project/6abf692aab636771b039c71a
% The measurement is the stellar linking length matched to an isophote
% in group-mass haloes. The Jacobi-radius finder paper is the separate file ms_jacobi.tex.
% Author order: Jaehyun Lee, Jihye Shin, Yonghwi Kim, Jaewon Yoo,
% M. James Jee, Jongwan Ko, Juhan Kim (corresponding).
% The group-mass sample is finished. The TNG300-1 merger-tree test uses
% 36 central branches common to three snapshots.
\PassOptionsToPackage{hyperfootnotes=false}{hyperref}
\documentclass[usenatbib]{mnras}
\usepackage{graphicx}
\usepackage{amsmath,amssymb}
\usepackage[T1]{fontenc}
\usepackage{newtxtext,newtxmath}
\usepackage{booktabs}
\usepackage{placeins}
\hyphenation{pixels}

\title[Stellar groups and surface-brightness boundaries]
{Matching stellar groups to surface-brightness boundaries: effects on inferred galaxy growth}

\author[Lee et al.]{%
Jaehyun Lee,$^{1}$
\newauthor
Jihye Shin,$^{1}$
Yonghwi Kim,$^{3}$
Jaewon Yoo,$^{1}$
M.~James Jee,$^{2}$
Jongwan Ko,$^{1}$
and Juhan Kim$^{4}$\thanks{E-mail: kjhan@kias.re.kr}
\\
$^{1}$Korea Astronomy and Space Science Institute, 776 Daedeok-daero, Yuseong-gu, Daejeon 34055, Republic of Korea\\
$^{2}$Department of Astronomy, Yonsei University, 50 Yonsei-ro, Seodaemun-gu, Seoul 03722, Republic of Korea\\
$^{3}$Korea Institute of Science and Technology Information, 245 Daehak-ro, Yuseong-gu, Daejeon 34141, Republic of Korea\\
$^{4}$Korea Institute for Advanced Study, 85 Hoegi-ro, Dongdaemun-gu, Seoul 02455, Republic of Korea
}

\begin{document}
\label{firstpage}
\pagerange{\pageref{firstpage}--\pageref{lastpage}}
\maketitle

\begin{abstract}
A surface-brightness limit defines the apparent extent of a galaxy in an image. We calibrate friends-of-friends (FoF) linking lengths against mock $V$- and SDSS $g$-band images at surface-brightness limits of $26.5$ and $27\,\mathrm{mag\,arcsec}^{-2}$. The images include cosmological dimming and redshifted spectra. Across 75 halo-output entries spanning six outputs of one Darwin volume, a 15-point grid using corrected luminosity distances gives a median best-fitting linking parameter of $0.09$ in both bands under either positional weighting or weights based on the fraction of each star's point-spread function (PSF) inside the isophote. In 130 CAMELS hosts, linking parameter $0.10$ raises mean full-image overlap relative to $0.20$ by $0.101/0.119$ in $V/g$; completeness falls by $0.040/0.039$.
Using a $0.5$-arcsec pixel-centre distance threshold around the isophote, median overlap is $0.516/0.528$ ($V/g$) for parameters from the original nine-point positional calibration and $0.575/0.592$ at $0.20$; at a $2$-arcsec threshold, the ranking reverses. In 29 TNG300-1 branches with central-pixel detections in every view at both endpoints, $g$-band isophotal stellar-mass growth exceeds central-subhalo growth by $0.196$ dex. Calibrated FoF groups have near-zero median growth bias, but a linking parameter of $0.20$ gives smaller mean absolute error. Pre-existing stars newly assigned to the central dominate the median central stellar-mass increase. At physical observer redshifts, particle-level age and metallicity weights reduce half-light-radius growth by $0.255$ dex in $V$ and $0.324$ dex in $g$ relative to uniform-weight controls. The photometric comparisons do not identify gravitationally unbound stars or the intracluster-light fraction.
\end{abstract}

\begin{keywords}
galaxies: groups: general -- galaxies: haloes -- methods: numerical
\end{keywords}

\section{Introduction}
\label{sec:intro}

Diffuse stellar light extends beyond individual galaxies in groups and
clusters and preserves evidence of stripping and mergers
\citep{zwicky1951,mihos2005,mihos2017,rudick2011,contini2021,montes2022}.
Observers often measure diffuse light in pixels fainter than a chosen
surface-brightness limit. A dynamical definition selects stars unbound
to every member galaxy. The stellar distribution has no sharp physical
edge, and the measured diffuse stellar mass depends on the definition.

A surface-brightness limit also marks an apparent boundary around a
central galaxy. The contour at the limit is an isophote. We call the
connected area brighter than the limit the isophotal region.
Limits near $26.5\,\mathrm{mag\,arcsec}^{-2}$ in $V$
and $27\,\mathrm{mag\,arcsec}^{-2}$ in $g'$ are widely used
\citep{feldmeier2004,zibetti2005,burke2015,kluge2020,montenegro2025}.
Profile fits, fixed apertures, and stacked images provide other ways
to measure diffuse light \citep{gonzalez2005,gonzalez2007,zhang2019,montenegro2023}.
Measurements from images do not test whether stars are gravitationally
bound. Spatial and dynamical definitions give different diffuse-light
fractions even in the same simulated clusters
\citep{rudick2011,cui2014,brough2024}.
Two classifications can give similar diffuse-light fractions but
different diffuse-light distributions around the central galaxy
\citep{montenegro2025}.
Sky subtraction, the point-spread function (PSF), cosmological dimming,
and redshifted spectra change the apparent galaxy boundary
\citep{tolman1930,tang2018}.
Comparing simulated stellar groups with images therefore requires an
explicit procedure for constructing the images.

An image contour also determines the stars included in a measured mass
and radius. Cosmological dimming can move the contour while a
simulated stellar distribution remains fixed. The inferred mass and size
growth can then combine physical evolution with image selection. We
separate the contributions by following central galaxies along merger
trees and rendering each stellar distribution at multiple observer
redshifts. Exact stellar-particle IDs separate changes in central-subhalo
membership from changes in particle masses, while controlled population
renderings measure how age and metallicity differences among stars,
including spatial gradients, change half-light-radius growth.

A friends-of-friends algorithm joins nearby particles through chains
of pairwise links without testing gravitational binding \citep{davis1985}.
PGalF identifies stellar density peaks and assigns self-bound galaxy
members \citep{kim2023}. We use the PGalF catalogue to remove resolved
satellites, then run a separate friends-of-friends calculation around
a stellar density peak. A dimensionless parameter $b_{\mathrm{link}}$ sets the linking
distance between star particles. Increasing $b_{\mathrm{link}}$ adds stars to the group.
The linking rule uses particle masses and the cosmic mean matter
density. The rule does not impose a single stellar-density boundary
\citep{more2011}.

The comparison asks how closely a three-dimensional stellar group can
reproduce the light distribution selected by a projected brightness limit.
We first calibrate the linking parameter against isophotes at
$26.5\,\mathrm{mag\,arcsec}^{-2}$ in $V$ and
$27\,\mathrm{mag\,arcsec}^{-2}$ in SDSS $g$ in one cosmological
hydrodynamical volume. The $g$ limit follows the $g'$ threshold of
\citet{montenegro2025}. We project each stellar group along three
perpendicular directions and compare the group light inside each
isophote with the combined light from group members and all stars inside
the isophote, counting each star once. We select the $b_{\mathrm{link}}$ giving the largest
light-weighted overlap.
We test fixed linking parameters on 49 CAMELS Swift-EAGLE groups from
boxes CV20--CV26, then evaluate the fixed values on a separate sample of
130 groups from boxes CV0--CV19. Neither sample is used to change the
Darwin calibration, and both use the same Swift-EAGLE model. We also test
40 TNG300-1 groups at one output with a fixed linking parameter and follow
36 central galaxies in that simulation along merger trees to measure
stellar-mass and size growth. Finally,
we render 75 Darwin entries from outputs at $z\leq1.990$ at five
observed redshifts while holding proper stellar positions and ages fixed.
The repeated renderings separate changes in the apparent galaxy boundary
from physical evolution of the stellar system. Additional tests vary image
construction and particle sampling and measure group members outside
the isophote and non-members inside. The comparison calibrates stellar
grouping against an apparent galaxy boundary and does not
measure an intracluster-light fraction.

\section{Simulation and sample}
\label{sec:sample}

\subsection{Darwin simulation}
\label{sec:volume}

We use star particles from the Darwin cosmological hydrodynamical
simulation. The periodic box is $44.171\,h^{-1}\mathrm{cMpc}$ on a side.
The flat cosmology has $\Omega_{m}=0.315$, $\Omega_{b}=0.045$, $\Omega_{\Lambda}=0.685$,
and $H_{0}=67.4\,\mathrm{km\,s}^{-1}\mathrm{Mpc}^{-1}$.
Darwin uses RAMSES \citep{teyssier2002} with mesh levels 10--17.
The finest cell spans
$0.337\,h^{-1}\mathrm{ckpc}$ at every output, corresponding to
$0.407$ proper kpc at $z=0.227$. Star particles have a minimum birth
mass of $10^{5}\,\mathrm{M}_{\odot}$. The supplementary methods give the
resolution conversion and star-formation, feedback, and black-hole
settings. For each halo, we fix the input list of stars before varying
the linking parameter. Stellar separations account for the periodic
box boundaries.
The supplementary methods record the input catalogue settings,
including disabled stellar linking at $z=0.227$.

\subsection{Halo selection}
\label{sec:cut}

We select haloes with friends-of-friends group masses of at least
$3\times10^{13}\,h^{-1}\mathrm{M}_{\odot}$ at $z\approx0.2$, and at least
$10^{13}\,h^{-1}\mathrm{M}_{\odot}$ at earlier outputs.
Both thresholds select massive galaxy groups. We refer to haloes above
$10^{14}\,h^{-1}\mathrm{M}_{\odot}$ as clusters and report halo masses using the
friends-of-friends group mass.
The six outputs at $z=3.002$, $1.990$, $1.002$, $0.500$,
$0.227$, and $0.190$ contain 1, 3, 17, 33, 10, and 12 haloes.
The 76 group masses span $1.050\times10^{13}$ to
$1.337\times10^{14}\,h^{-1}\mathrm{M}_{\odot}$, with three haloes above
$10^{14}\,h^{-1}\mathrm{M}_{\odot}$. One halo has a maximum at a grid endpoint
and is excluded, leaving 75 halo-output entries for the calibration.
Nine entries at $z=0.227$ have positional matches at $z=0.190$.
Exact dark-matter particle IDs establish links through the $z=0.500$
output. Links from $z=0.500$ to the two lower-redshift outputs use
positions and remain provisional.
We therefore treat 75 as the number of halo-output measurements, not
as the number of independent systems. The supplementary material lists
every catalogue entry and the corresponding sample-selection flags.
The $z=0.227$ output uses a different stellar-selection configuration, so
we test prediction within the $z=0.1904$ output using the 12 retained
haloes. For each test, we omit one halo, take the median of the three
projection-specific $b_{\mathrm{link}}$ values giving maximum overlap
for each of the other 11 haloes, and then take the median across those
11 calibration haloes. We evaluate the omitted halo in all three
projections. We report the $z=0.1904$ test separately from the six-output
calibration.

\subsection{Previously analysed CAMELS Swift-EAGLE sample}
\label{sec:camels}

We test the fixed $b_{\mathrm{link}}=0.10$ calibration on 49 groups from the first-generation
CAMELS Swift-EAGLE L25n256 CV boxes 20--26 \citep{villaescusa2021}.
We use the $z=0$ snapshot 090 and the corresponding Groups catalogue
(\texttt{snapshot\_090.hdf5} and \texttt{groups\_090.hdf5}). The box is
$25\,h^{-1}\mathrm{cMpc}$ on a side and contains $256^3$ initial gas
elements and dark-matter particles. The snapshot header gives an initial gas
element mass of $1.89\times10^7\,\mathrm{M}_{\odot}$. Stellar particle masses
vary after stellar mass loss. The catalogue parameters give a maximum
proper stellar softening of $5.84\,h^{-1}\mathrm{pkpc}$. The selected
friends-of-friends masses span $1.00$--$9.14\times10^{13}\,h^{-1}\mathrm{M}_{\odot}$.
We use catalogue particle IDs and the CAMELS cosmology
($\Omega_m=0.3$, $h=0.6711$). We render the systems at
$z_{\rm obs}=0.1904$ with the same filters, surface-brightness limits,
pixel scale, and point-spread function as the Darwin images. We compare
fixed $b_{\mathrm{link}}=0.10$ and $b_{\mathrm{link}}=0.20$ without recalibrating $b_{\mathrm{link}}$ on the CAMELS sample.
The central galaxy is the subhalo with the most star particles. Resolved
satellites with at least 100 star particles are masked. For each selected
stellar set, projected $R_{50}$ encloses half of the total stellar flux
of the selected set. The aperture is a two-dimensional projected circle.
The image PSF determines the isophotal region, while luminosities and
$R_{50}$ use full particle fluxes at projected positions, without PSF
aperture weights. The comparison uses one finite-resolution simulation
and does not establish convergence with particle or force resolution. We use the central
subhalo's minimum-potential position as the common centre for the image,
the 30 pkpc aperture, and projected radii. The CAMELS and Darwin host-mass
ranges overlap, while Darwin extends to higher masses. The test uses a
simulation-known centre and does not test automatic centre identification
from an observed image.

\subsection{CV0--CV19 validation sample}
\label{sec:camels-validation}

We select 130 FoF hosts from the first-generation CAMELS Swift-EAGLE
L25n256 boxes CV0--CV19 at snapshot 090. The boxes do not overlap the
CV20--CV26 sample. The selected FoF masses span
$10^{13}$--$10^{14}\,h^{-1}\mathrm{M}_{\odot}$, and each central subhalo
contains at least 100 star particles. We render the $z=0$ systems at
$z_{\rm obs}=0.19044$ with the same filters, surface-brightness limits,
pixel scale, PSF, satellite masks, and stellar-flux calculation as the
CV20--CV26 sample. Catalogue particle IDs define each group. The central
subhalo's potential minimum, \texttt{SubhaloPos}, centres the image,
aperture, and projected radii. We compare $b_{\mathrm{link}}=0.20$,
$b_{\mathrm{link}}=0.10$, the low-redshift Darwin values of 0.12 in $V$
and 0.10 in $g$, and a projected 30 pkpc aperture. We do not tune
the linking parameter using the CV0--CV19 measurements. The centre rule
was selected during earlier work, before the present validation
measurements; the full analysis protocol was therefore not set in advance.
The validation measures transfer across boxes in one simulation model, not transfer
across simulation models. We reduce the three projections to one median
per host, and the 130 hosts are the statistical units. Bootstrap intervals resample boxes first and
hosts within each sampled box second.
The 30 pkpc projected aperture is a fixed baseline, not a radius fitted
on the Darwin calibration sample. The comparison therefore does not establish
FoF superiority over an optimized projected aperture. In the separate
TNG growth test, 3D spherical-aperture radii are selected from the 12
Darwin groups at $z=0.190$ and held fixed in TNG; that comparison uses a
different aperture geometry and does not calibrate a projected radius.

For the CAMELS validation sample, we mark a target contour as invalid
if the nearest pixel above the threshold lies in a separate isophotal
component rather than around the central galaxy, if the selected region
is empty or has zero luminosity, or if it touches the image edge. In the TNG
transfer test, a contour is invalid if the selected region is missing or
touches the image edge. We assign zero IoU and completeness to an invalid
view; precision and luminosity and $R_{50}$ errors are undefined. A host
counts once as a large failure if any target contour is invalid, median
contamination exceeds 0.20, or either host-median logarithmic ratio has
an absolute value greater than $\log_{10}2$. The ratios are
$\log_{10}(L_{\rm selected}/L_{\rm contour})$ and
$\log_{10}(R_{50,\rm selected}/R_{50,\rm contour})$. We report a separate
orientation-level fraction that counts a host when any one view crosses
the same threshold. The 90th-percentile
absolute luminosity and $R_{50}$ errors use valid contours only.

\subsection{TNG300-1 sample}
\label{sec:tngsample}

We select 40 friends-of-friends (FoF) hosts from the TNG300-1 simulation at
snapshot 84 ($z=0.1973$), before calculating any image measurements. The sample
covers FoF masses from $3.09\times10^{13}$ to
$1.288\times10^{14}\,h^{-1}\mathrm{M}_{\odot}$, with a median of
$4.680\times10^{13}\,h^{-1}\mathrm{M}_{\odot}$, and matches the Darwin host-mass
range. TNG300-1 is part of the IllustrisTNG simulation suite
\citep{nelson2019tng,pillepich2018tng,springel2018tng,nelson2018tng,naiman2018tng,marinacci2018tng}.
The dark-matter particle mass is $5.9\times10^{7}\,\mathrm{M}_{\odot}$,
the target baryonic mass resolution is $1.1\times10^{7}\,\mathrm{M}_{\odot}$,
and the dark-matter and stellar gravitational softening is 1.48 proper kpc
at $z\leq1$ \citep{nelson2018tng}. The tree analysis therefore samples
extended stellar structure at finite particle and force resolution; we do
not test numerical convergence.
For each host, we select the qualifying subhalo with the largest number of
stellar particles. We centre the images and projected apertures on the
subhalo's \texttt{SubhaloPos}, defined by the maximally bound particle
position and commonly used as the galaxy centre. The mass-weighted
\texttt{SubhaloCM} includes all bound particle and cell types and can be
displaced by asymmetric or distant material.\footnote{The field definitions
and units appear in the IllustrisTNG data specifications at
\url{https://www.tng-project.org/data/docs/specifications/}. TNG position
guidance recommends \texttt{SubhaloPos} for galaxy centring at
\url{https://www.tng-project.org/data/forum/topic/295/how-to-get-galaxy-positions/}.}
The stellar-density peak remains the root for linked stellar-group
selection. We retain the same satellite mask, filters, surface-brightness
thresholds, pixel scale, PSF, and flux calculation used for Darwin. We
apply the low-redshift Darwin calibration (0.12 in $V$ and 0.10 in $g$)
without adjustment. Fixed linking parameters of 0.10 and 0.20 and a
projected 30 pkpc circular aperture provide comparisons. We reduce the
three orthogonal views to one median per host. The 40 host haloes are the
statistical units. Earlier TNG overlap measurements using another image-centre rule had
already been inspected. The present centre check is post hoc, and no
linking parameter was tuned using TNG measurements.

For the growth analysis, we retrieve the SubLink main-progenitor branch
for each of the 40 descendants \citep{rodriguezgomez2015}. We retain the
36 branches whose central subhalo remains the FoF central at snapshots
84, 67, and 50 ($z=0.197$, 0.503, and 0.997). The catalogue-based cut
does not depend on image measurements. At each output, exact central-subhalo star
membership defines the bound-stellar reference. We render each physical
state at the simulation redshift and again at fixed observer redshift
$z_{\rm obs}=0.197$. The fixed-redshift control images isolate the
effect of changing observer conditions. We compare the
central-subhalo reference with isophotal selection, FoF groups with
$b_{\mathrm{link}}=0.10$, 0.12, and 0.20, a projected 30 pkpc circle, and
spherical apertures of radius 80 pkpc in $V$ and 60 pkpc in $g$. The
$b_{\mathrm{link}}=0.12$ ($V$) and 0.10 ($g$) values are the low-redshift
Darwin calibrations; 0.10 in $V$ and 0.20 in both bands are comparison
settings. The sphere radii were selected on the 12 Darwin groups at
$z=0.190$ from a grid of 10, 15, 20, 30, 45, 60, 80, and 120 pkpc by
maximizing the median host-level, three-view light-weighted overlap. The
full-sample radii are held fixed for TNG and are not refit using TNG
images. For every method, we calculate growth differences in each
projection, take the median within each branch, and bootstrap across
branches. We also calculate projected half-mass radii using current
stellar-particle masses in place of luminosity weights.

For each rendered image, we record whether the central pixel lies inside
the isophotal region, whether the selection uses the region around the
nearest pixel brighter than the threshold, the selected region's area,
and contact with the image edge. The nearest-pixel selection is
retained and flagged separately from a central-pixel detection. An image
with no above-threshold pixel has an empty selection.

\subsection{Artificial redshift experiment}
\label{sec:artificialredshift}

To measure changes caused by image construction, we render 75 halo-output
entries from outputs at $z\leq1.990$ at $z_{\rm obs}=0.190$, $0.227$, $0.500$, $1.002$, and $1.990$.
We hold the proper stellar positions, ages, and masses fixed at each
simulation output and convert proper positions to angles using the
angular-diameter distance at each observed redshift. We compare the full
image including dimming and redshifted spectra with control images. The
controls remove the $(1+z)^4$ surface-brightness
dimming, replace the redshifted stellar spectrum with the $z=0$ spectrum
while retaining luminosity-distance dimming, or use the $z=0$ spectrum
without luminosity-distance dimming. The final case is an extreme
no-distance control, not just the Tolman-factor correction. We keep the
angular pixel scale and point-spread function fixed. The physical
resolution therefore changes with $z_{\rm obs}$. Across rendering redshifts,
cosmological dimming, sampled wavelength, and physical resolution change
together. The renderings hold halo structure fixed and do not represent
galaxy evolution.
We omit sky noise, dust attenuation, and absorption by intergalactic gas.

\section{Mock images}
\label{sec:image}

\subsection{Stellar light}
\label{sec:light}

We calculate stellar light using python-fsps simple stellar populations
\citep{conroy2009}, a Chabrier initial mass function, and the FSPS
filters \texttt{v} and \texttt{sdss\_g}, with no interstellar dust
attenuation or nebular emission.
Absorption by intergalactic gas is also omitted.
We use AB magnitudes in both filters.
The population grid contains 32 ages from $0.01$ to
$13.5\,\mathrm{Gyr}$ and 12 metallicities from $10^{-4}$ to $0.05$,
with $Z_{\odot}=0.0142$ in FSPS.
We interpolate log stellar flux bilinearly in log age and log metallicity.
Ages and metallicities outside the grid are set to the nearest grid limit.
For a particle with formed mass $M_{\mathrm{form}}$ in solar masses,
the relative AB flux is
\begin{equation}
f=M_{\mathrm{form}}\,10^{-0.4m_{\mathrm{SSP}}(z)}.
\label{eq:stellarflux}
\end{equation}
The simple stellar population (SSP) magnitude includes the redshifted
spectrum and the FSPS luminosity-distance normalization.
The stored formed mass is divided by $h$ before multiplication.
Surface brightness is $-2.5\log_{10}(f_{\mathrm{pixel}}/A_{\mathrm{pixel}})$,
with pixel area in square arcseconds.
We calculate angular sizes and luminosity distances using the simulation
cosmology. The archived FSPS grids use a different cosmology for
luminosity distances. We multiply the archived stellar fluxes by
$(D_{L,\mathrm{FSPS}}/D_{L,\mathrm{sim}})^2$ before generating the main
images. The archived normalization is used in the explicitly labelled
illustrative images and sensitivity diagnostics. The primary numerical
results use the simulation cosmology.
The archived flux grids accompany the supplementary tables.
Recomputing every grid value with python-fsps 0.4.7 reproduces the
archived fluxes exactly. The original FSPS distance calculation uses
$\Omega_m=0.27$ and $H_0=72\,\mathrm{km\,s}^{-1}\mathrm{Mpc}^{-1}$.
The distance correction makes the images $0.046$--$0.130$ magnitudes
fainter. The supplementary methods give the distance factors and compare
the corrected and archived normalizations.

We select the catalogue galaxy containing the largest number of star
particles as the central galaxy. The selected galaxy has the largest
stellar mass in all 76 haloes and the largest luminosity in 52 haloes
in $V$ and 52 in $g$. We remove every other catalogue galaxy
represented by at least $100$ star particles from the images and linking
calculation, using the same satellite list for every $b_{\mathrm{link}}$.
Catalogue galaxies with fewer than $100$ stars and background stars remain.
The image centre is the mass-weighted centre of the dark matter, the gas,
the sink particles, and the stars assigned to the selected central galaxy.

\subsection{Point-spread function}
\label{sec:psf}

We convolve each image once with the same PSF for every $b_{\mathrm{link}}$ and view.
The core is a circular Moffat profile with
$\mathrm{FWHM}=1.0\,\mathrm{arcsec}$ and $\beta=4.765$.
An $A/r^{2}$ wing is added to the Moffat core.
The DECam amplitudes are $A_{g}=4.5\times10^{-4}$ and
$A_{r}=3.3\times10^{-4}$ for a unit-flux star.
The kernel extends to $80\,\mathrm{arcsec}$ and has unit total flux.
Since DECam has no $V$ filter, we use the $r$-band amplitude for the
$V$ wing and treat $g$ as the primary band.
The supplementary methods give the wing transition between $7$ and
$8\,\mathrm{arcsec}$.
The pixel scale is $0.2\,\mathrm{arcsec}$.
We add each star's flux to the nearest pixel within the image.
The image extends $180\,\mathrm{arcsec}$ either side of the centre and
contains $1800$ pixels on a side.
Stars outside the image contribute no flux. Convolution preserves the
image dimensions and discards light spread beyond the image edges.
We treat pixels sharing a side or corner as connected.
The images contain no sky noise.
Section~\ref{sec:sensitivity} reports tests with added sky residuals.

\label{sec:axes}

We project each halo along the three box axes using the same pixel
scale, PSF, centre, and satellite mask.

\section{Matching a linking length to an isophote}
\label{sec:match}

\subsection{Galaxy catalogue and stellar linking}
\label{sec:finder}

PGalF separates stellar density peaks with a watershed algorithm,
merges peaks with prominence below $0.30$, and discards regions with
fewer than $100$ particles \citep{kim2023}.
PGalF adds members through density shells, truncates satellites at the
tidal radius, and removes unbound stars around satellite cores.
The catalogue fixes the satellite mask and stellar input list for the
separate friends-of-friends calculation.
The supplementary methods define peak prominence, describe the
watershed and satellite binding calculation, and illustrate the
galaxy boundaries.

\label{sec:group}

The dimensionless friends-of-friends parameter $b_{\mathrm{link}}$ sets the linking
distance $\ell$ through
\begin{equation}
\ell
=
b_{\mathrm{link}}\left(\frac{m_{\star}}{\Omega_{m}\rho_{c}}\right)^{1/3},
\label{eq:link}
\end{equation}
with $\rho_{c}=2.7755\times10^{11}\,h^{2}\,\mathrm{M}_{\odot}\,\mathrm{Mpc}^{-3}$.
Here $m_{\star}$ is the particle's current stellar mass. The separately
stored formed mass is used only to calculate flux, after division by
$h$. In catalogue units, the factor $h^{2}$ in $\rho_{c}$ cancels, giving
\begin{equation}
\ell
=
b_{\mathrm{link}}\left(\frac{m_{\star}}{2.7755\times10^{11}\,\Omega_{m}}\right)^{1/3},
\label{eq:linkcode}
\end{equation}
with $m_{\star}$ in $h^{-1}\mathrm{M}_{\odot}$ and $\ell$ in $h^{-1}$ comoving Mpc.
Two stars are linked when the comoving separation is smaller than the
mean of the two individual linking distances.
We link only star particles. The normalization of $b_{\mathrm{link}}$ uses the cosmic
mean matter density \citep{davis1985} and does not use the mean stellar
separation.
A particle of $10^{5}\,h^{-1}\mathrm{M}_{\odot}$ has
$\ell=1.046\,h^{-1}\mathrm{ckpc}$ at $b_{\mathrm{link}}=0.10$.
The minimum birth mass is $10^{5}\,\mathrm{M}_{\odot}$, and stellar particle
masses vary within a halo. The linking rule does not follow a single
density contour. The supplementary methods discuss equal-mass density
estimates and the percolation interpretation of \citet{more2011}.
For each star, we find the radius containing the star and 31 neighbours
from the input list. We start the group at the star with the smallest
radius. The starting star, called the
density seed, lies in a dense stellar region but need not lie near the
image centre. The seed is separate from the watershed peak used by
PGalF to identify the galaxy in Section~\ref{sec:finder}.
For images including cosmological dimming and redshifted spectra, we
test $b_{\mathrm{link}}=0.02$, $0.03$, $0.05$, $0.08$,
$0.10$, $0.12$, $0.15$, $0.18$, and $0.20$.
We exclude a halo from the calibration if any view gives the largest
overlap at the first or last tested $b_{\mathrm{link}}$. The best match may then lie
beyond the tested range.

\subsection{Light-weighted overlap}
\label{sec:iou}

We select a connected region of pixels brighter than
$26.5\,\mathrm{mag\,arcsec}^{-2}$ in $V$ or
$27\,\mathrm{mag\,arcsec}^{-2}$ in $g$.
If the centre pixel is brighter than the limit, we use the connected
region containing the centre. Otherwise, we use the region
containing the nearest pixel brighter than the limit. We select no
region when every pixel is fainter than the limit.
The linking calculation selects a three-dimensional stellar group.
The isophote selects stars by projected position from the same input
list. We divide the light
from group members inside the isophote by the combined light from the
group and stars inside the isophote, counting each star once.
The light-weighted intersection over union (IoU) is defined as
\begin{equation}
\mathrm{IoU}
=
\frac{L_{\mathrm{both}}}
{L_{\mathrm{group}}+L_{\mathrm{inside}}-L_{\mathrm{both}}}.
\label{eq:iou}
\end{equation}
Here $L_{\mathrm{group}}$ is the total light of group members,
$L_{\mathrm{inside}}$ is the total light of stars projected inside the
isophote, and $L_{\mathrm{both}}$ is the light of group members projected
inside the isophote. Subtracting $L_{\mathrm{both}}$ in the denominator
prevents double counting. Each star contributes the full stellar flux
according to the projected position relative to the isophote in the
convolved image. Thus the PSF locates the contour but does not redistribute
the flux used in the primary IoU, selected luminosity, or $R_{50}$. We report
a separate calculation weighting each star by the fraction of the stellar PSF
inside the contour below.
Completeness, $L_{\mathrm{both}}/L_{\mathrm{inside}}$, measures how much
stellar light inside the isophote belongs to the group. Precision,
$L_{\mathrm{both}}/L_{\mathrm{group}}$, measures how much group light lies
inside the isophote. Neither fraction tests gravitational binding.
For each view, we choose the tested $b_{\mathrm{link}}$ giving the highest overlap and
choose the smallest $b_{\mathrm{link}}$ in a tie. For each halo and band, we take the
median of the three selected values after applying the endpoint
exclusion described in Section~\ref{sec:group}.

We also assign flux using the fraction $q_i$ of star $i$'s normalized
PSF inside the same isophote. For the set of group members $G$,
the fluxes in Equation~\eqref{eq:iou} become
\begin{equation}
L_{\mathrm{group}}=\sum_{i\in G}f_i,\qquad
L_{\mathrm{inside}}=\sum_i q_i f_i,\qquad
L_{\mathrm{both}}=\sum_{i\in G}q_i f_i.
\label{eq:psfscore}
\end{equation}
The group total retains every member's full flux.
Convolution can spread light from a star outside the isophote into the
isophotal region. Stars outside the image have $q_i=0$, consistent
with image construction. Changing the flux weights leaves stellar
group membership unchanged.
Completeness and precision use the same fluxes in
$L_{\mathrm{both}}/L_{\mathrm{inside}}$ and
$L_{\mathrm{both}}/L_{\mathrm{group}}$.

\label{sec:noticl}

Stars near the image centre can dominate Equation~\eqref{eq:iou}.
We therefore also examine group members outside the isophote and
non-members inside, particularly near the boundary.

\section{Results}
\label{sec:results}

\subsection{Within-output validation at redshift 0.190}
\label{sec:withinoutput}

Every test that omits one halo selects $b_{\mathrm{link}}=0.12$ in $V$
and $0.10$ in $g$. The median IoU of the omitted haloes is 0.959 in $V$
and 0.962 in $g$, compared with 0.927 and 0.921 at $b_{\mathrm{link}}=0.20$.
All 12 haloes have higher median IoU at the selected parameter. Median
completeness is 0.983/0.980 at the selected parameters and 0.995/0.996
at $b_{\mathrm{link}}=0.20$ ($V/g$). Median precision rises from
0.932/0.925 at $b_{\mathrm{link}}=0.20$ to 0.975/0.985 at the selected
parameters. Median IoU falls below 0.90 for 3/12 $V$ haloes and 2/12 $g$
haloes. At least one projection falls below 0.90 for 3/12 haloes in each
band. The supplementary table gives results for every omitted halo and
projection. The resampling tests share calibration haloes and use one simulation
output, so they describe within-output prediction but do not provide an
independent-volume error estimate.

\subsection{Calibration across the six outputs}
\label{sec:bz}

All three views have maximum overlap between the grid endpoints for
75 halo measurements. With luminosity distances evaluated in the Darwin
cosmology, the median selected $b_{\mathrm{link}}$ is $0.08$ in $V$ and $0.10$ in $g$. The 16th--84th percentile range is $0.08$--$0.12$ in both bands. The
median overlap is $0.9658$ in $V$ and $0.9721$ in $g$.
Table~\ref{tab:bmed} gives the results at each output, and
Figure~\ref{fig:bmatch} shows the median selected $b_{\mathrm{link}}$ for each halo.
At $z=0.190$--$0.227$, the medians are $0.12$ in $V$ and $0.10$ in $g$.
One halo at $z=0.227$ has maximum overlap at a grid endpoint and is
excluded from the calibration (Section~\ref{sec:fail}).
The mass cut changes between outputs, nine systems recur at the last
two outputs, and only one halo enters the sample at $z=3.002$. All entries
come from one simulation volume, so output-to-output differences do not
measure redshift evolution.
An equal-output sensitivity averages the six output-specific median IoU
curves and selects $b_{\mathrm{link}}=0.10$ in both bands. The entry-level
medians remain $0.08$ in $V$ and $0.10$ in $g$. The $V$-band shift reflects
the joint change in output weighting and estimator: the primary summary is
the median of entry-level best parameters, whereas this sensitivity selects
the maximum of an equally weighted mean of output-specific median-IoU
curves. The shift cannot be attributed to output weighting alone. The
primary summary targets a typical retained halo-output entry under the
stated mass cuts. For each view, the loss is one minus the full-image,
light-weighted IoU; we take the median of the three view-specific best
parameters for each entry, then the median across entries. The six outputs
share one volume and include repeated systems, so the sensitivity is
descriptive and has no independent-output uncertainty.
Correlations between halo mass and selected $b_{\mathrm{link}}$ change sign between
outputs and show no consistent mass dependence.

Using the archived FSPS luminosity-distance normalization instead gives
median parameters of $0.10$ in both bands. The distance correction
changes the selected parameter for 8 haloes in $V$ and 11 in $g$. The
paired median parameter change is zero in both bands. The corrected
values are the main results because the calculation uses one cosmology for both
angular and luminosity distances.

For each halo, we subtract the smallest selected $b_{\mathrm{link}}$ among the three
views from the largest. We also measure the interval of tested $b_{\mathrm{link}}$
values giving overlap within $0.02$ of the maximum in each view, then
take the median interval width over the three views.
Across the 75 haloes, the median range across views is $0.02$ in both
bands. The median interval width is $0.05$ in both bands.
Several nearby values of $b_{\mathrm{link}}$ therefore give similar light ratios.
The intervals describe the breadth of the overlap curves, not
statistical uncertainties, and can extend beyond the tested grid.

\begin{table}
\centering
\caption{Linking parameters giving maximum light-weighted overlap in
images with cosmological dimming and redshifted spectra. Fluxes and
angular sizes use the Darwin cosmology.
$N$ counts haloes with maximum overlap between the first and last tested
values of $b_{\mathrm{link}}$ in all three views.
At $z=0.227$ the output contains ten haloes, and one is omitted
(Section~\ref{sec:fail}).
$b_{\mathrm{link}}$ is the median across haloes after taking the median over views for
each halo and is selected from the original nine-point grid.
$\mathrm{IoU}$ is the median across haloes of the median overlap over
views, with $b_{\mathrm{link}}$ selected separately for each view.}
\label{tab:bmed}
\small
\setlength{\tabcolsep}{2.5pt}
\begin{tabular}{lrrrrr}
\toprule
$z$ & $N$ & $b_{V}$ & $b_{g}$ & $\mathrm{IoU}_{V}$ & $\mathrm{IoU}_{g}$ \\
\midrule
3.002 & 1 & 0.10 & 0.12 & 0.9188 & 0.9396 \\
1.990 & 3 & 0.10 & 0.15 & 0.9525 & 0.9546 \\
1.002 & 17 & 0.08 & 0.08 & 0.9600 & 0.9725 \\
0.500 & 33 & 0.08 & 0.08 & 0.9696 & 0.9754 \\
0.227 & 9 & 0.12 & 0.10 & 0.9726 & 0.9763 \\
0.190 & 12 & 0.12 & 0.10 & 0.9669 & 0.9664 \\
\bottomrule
\end{tabular}
\end{table}

\begin{figure}
\centering
\includegraphics[width=\columnwidth]{figures/b_match.pdf}
\caption{Linking parameter $b_{\mathrm{link}}$ giving the largest light-weighted overlap
between the projected distribution of group members and the distribution
of stars inside the isophote.
Each point gives the median of
three viewing directions for one halo. The horizontal axis gives simulation
output redshift $z$, with outputs spaced equally. Small horizontal offsets
separate haloes at the same output. Filled circles mark median overlap
at or above $0.90$, and open squares mark median overlap below $0.90$.
The cross marks the halo with maximum overlap at a tested endpoint; the
halo is displayed but excluded from calibration.}
% Alt text: Selected linking parameter versus simulation redshift for 75 halo-output entries. The plotted value for each halo is the median from three viewing directions. Symbols distinguish overlaps above or below 0.90, and the cross marks an endpoint maximum excluded from calibration.
\label{fig:bmatch}
\end{figure}

A representative cluster image using the archived photometric normalization is retained as Supplementary Fig. S1.

\subsection{Comparison with the previously adopted linking parameter}
\label{sec:comparison}

$b_{\mathrm{link}}=0.20$ is the legacy setting in the stellar FoF reader
and parent catalogues. The setting was not calibrated against an isophotal
boundary and does not define a universal stellar standard. We rerun FoF at
$b_{\mathrm{link}}=0.10$ and $0.20$ on the same input star list for each
halo.

Using $b_{\mathrm{link}}=0.10$ for every halo gives median overlaps of $0.9582$ in $V$
and $0.9626$ in $g$, after taking the median over three views for each
halo. The previously adopted $b_{\mathrm{link}}=0.20$ gives $0.9198$ and $0.9328$.
For each halo, we subtract the median overlap at $b_{\mathrm{link}}=0.20$ from the
median at $b_{\mathrm{link}}=0.10$. The median increases are $0.0363$ and $0.0339$.
Overlap increases for 72 of the 75 halo-output entries in $V$ and 68 in
$g$, and on 209 of the 225 views in $V$ and 202 in $g$. Repeated outputs
are not independent systems. The counts describe the sample and do not
provide independent-sample confidence intervals.

Reducing $b_{\mathrm{link}}$ excludes group light outside the isophote but also omits
more stellar light inside. At $b_{\mathrm{link}}=0.10$, median completeness and precision
across halo summaries are $(0.9840,0.9812)$ in $V$ and $(0.9866,0.9772)$
in $g$. At $b_{\mathrm{link}}=0.20$, completeness increases to $0.9975$ and $0.9981$,
while precision falls to $0.9211$ and $0.9341$.
Figure~\ref{fig:bcurves} shows the ratios as functions of $b_{\mathrm{link}}$, and
Figure~\ref{fig:bresidual} shows completeness and precision at each
view's selected $b_{\mathrm{link}}$.
Each coloured curve takes the median overlap across haloes at each
tested $b_{\mathrm{link}}$. Table~\ref{tab:bmed} instead takes the median of the $b_{\mathrm{link}}$
values selected separately for individual haloes. The two operations
can give different $b_{\mathrm{link}}$ values when several adjacent grid points have
similar overlap.

The linking parameter also changes the total stellar light assigned
to the central group. Relative to the light from stars inside the
isophote, the median absolute luminosity difference falls from
$8.19$ to $2.36$ per cent in $V$ and from $6.91$ to $1.98$ per cent in
$g$ when $b_{\mathrm{link}}$ decreases from $0.20$ to $0.10$.
We take absolute differences before the median over views and then
the median over the 75 haloes. The difference decreases for 69 haloes
in $V$ and 67 in $g$. Supplementary archived-normalization tables give exploratory radial light sums
and signed differences. The exploratory values do not enter the corrected primary results. The luminosities use stellar positions and
full stellar fluxes, as in equation~\ref{eq:iou}. The comparison shares
the light sums used for calibration and does not independently test
gravitational membership.

\begin{figure*}
\centering
\includegraphics[width=\textwidth]{figures/b_curves.pdf}
\caption{Overlap and completeness versus linking parameter.
Columns show $V$ and SDSS $g$. Rows show $\widetilde{\mathrm{IoU}}$ and
$\widetilde{C}$, where $C$ is completeness and tildes denote the median
over three projections per entry followed by the median across entries.
Curves give these medians at each tested linking parameter.
All panels show y values from $0.70$ to $1.00$. Values below $0.70$ fall
outside the plotted range. Only entries with an interior maximum in both
bands are included.
The overlap uses the corrected luminosity-distance convention described
in the text.}
% Alt text: Four panels show median overlap and completeness versus linking parameter in V and g. Overlap peaks near the selected parameter, while completeness remains high and generally increases toward larger linking lengths.
\label{fig:bcurves}
\end{figure*}

\begin{figure*}
\centering
\includegraphics[width=\textwidth]{figures/b_residual.pdf}
\caption{Fractions of stellar light inside the isophote belonging to the
group (completeness) and group light inside the isophote (precision),
measured at each view's highest overlap on the original grid. Each point is one view, and the diagonal marks equal completeness and
precision. Filled circles have overlap at least $0.90$, open squares have
interior maxima below $0.90$, and crosses mark maxima at a grid endpoint.
Both panels use the same $0.60$--$1.00$ range on each axis. Points
below $0.60$ fall outside the displayed range. Colour denotes simulation
output redshift.}
% Alt text: Two scatter plots compare completeness with precision at each view's selected linking parameter, one for V and one for g. The diagonal marks equal values; marker shapes identify high-overlap cases, low-overlap cases, and grid-edge maxima, while colour marks redshift.
\label{fig:bresidual}
\end{figure*}

\subsection{View dependence and calibration checks}
\label{sec:adopted}

Applying each halo's median selected $b_{\mathrm{link}}$ to all three views gives median
overlaps of $0.9645$ in $V$ and $0.9686$ in $g$, compared with
$0.9658$ and $0.9721$ when $b_{\mathrm{link}}$ is selected separately for each view.
The selected value differs between views for 49 of the 75 entries in
both bands. The median loss is zero. The largest loss is $0.030$ in $V$
and $0.068$ in $g$. The calibration uses the original nine-point grid.
On the corrected-distance 15-point grid, the median of the individual
best parameters is $0.09$ in both bands for both positional and
PSF-aperture weights. The median-over-entry positional-overlap curves peak
at $b_{\mathrm{link}}=0.09$ in both bands; the PSF-aperture curves peak
at $0.09$ in $V$ and $0.10$ in $g$. At the individual PSF-weighted
optima, median overlap is $0.946$ in $V$ and $0.951$ in $g$, compared
with $0.968$ and $0.973$ when each star contributes the full stellar flux based
on position. At fixed $b_{\mathrm{link}}=0.10$, PSF-weighted median
overlap is $0.937$ and $0.944$, compared with $0.905$ and $0.917$ at
$0.20$. The two weighting rules therefore give similar linking
parameters but different overlap values; the Supplement defines both
objectives and reports the full curves.

The corrected image centre pixel lies inside the isophote in 186 of 225
$V$ projections and 189 of 225 $g$ projections. The remaining 39 and 36
projections use the nearest above-threshold region. No selected region is
empty or touches the image edge. Requiring a centre-pixel detection in all
three projections leaves 55 $V$ and 53 $g$ halo entries. Their median
coarse-grid parameters remain $0.08$ in $V$ and $0.10$ in $g$; the
fine-grid medians are $0.09$ for positional and PSF-aperture weighting
in both bands. The Supplement reports scores for this common-detection
subset. The
separate CV0--CV19 validation already requires a detected central component
and has no invalid target among the 130 hosts in either band. In the
TNG300-1 growth analysis, applying the same endpoint rule to every
selection retains 36 branches in $V$ and 29 in $g$.

We first select a common parameter while omitting each output and nearby
haloes from the calibration sample; each test selects $b_{\mathrm{link}}=0.10$.
Because these tests can share repeated systems, we also omit all entries
in groups connected by exact dark-matter particle IDs. Among 53 entries
at $z=1.990$, $1.002$, and $0.500$, this gives 33 groups; all 33 tests
select $b_{\mathrm{link}}=0.10$ in both bands. The median omitted overlap is $0.956$ in $V$ and $0.962$ in $g$.
The equal-group mean gain over $b_{\mathrm{link}}=0.20$ is $0.049$ ($0.042$--$0.055$)
in $V$ and $0.042$ ($0.034$--$0.049$) in $g$ (16th--84th percentile
bootstrap intervals). The particle-ID check covers only outputs with exact particle-ID
links and does not establish a complete set of independent halo lineages.

\subsection{Transfer to CAMELS Swift-EAGLE}
\label{sec:camelsresults}

At fixed $b_{\mathrm{link}}=0.10$, the equal-host mean of the three-view median overlap
is $0.894$ in $V$ and $0.872$ in $g$. Relative to $b_{\mathrm{link}}=0.20$, the paired
mean overlap increases by $0.094$ in $V$ (95 per cent box-then-host
bootstrap interval $0.083$--$0.106$) and $0.111$ in $g$
($0.098$--$0.127$). Median overlap remains below $0.90$ for 65 per cent
of hosts in $V$ and 82 per cent in $g$.
Reducing $b_{\mathrm{link}}$ from $0.20$ to $0.10$ lowers the mean fraction of selected
light outside the isophote by $0.131$ in $V$ and $0.145$ in $g$,
while completeness decreases by $0.035$ in both bands.

At $b_{\mathrm{link}}=0.10$, mean logarithmic selected-to-isophotal ratios are
$+0.006/+0.017$ dex for luminosity and $+0.012/+0.029$ dex for projected
half-light radius ($V/g$). The 90th-percentile absolute errors are
$0.033/0.051$ dex in luminosity and $0.048/0.081$ dex in projected
half-light radius.
Absolute errors use $|\log_{10}(X_{\rm selected}/X_{\rm isophote})|$,
with each host represented by the median over three projections.
The fixed 30 pkpc projected circular aperture gives larger
90th-percentile luminosity errors of $0.166/0.158$ dex and radius errors
of $0.238/0.236$ dex ($V/g$). The aperture selects less light outside
the isophote but misses more light inside. Figure~\ref{fig:transfer}
compares host-level overlap, luminosity error, and half-light-radius
error for the two FoF choices and the aperture. In the direct paired
comparison, the mean overlap gain over the aperture is $+0.040$
($+0.005$--$+0.083$) in $V$ and $+0.010$ ($-0.024$--$+0.051$) in $g$.
The $g$ interval includes zero. The mean outside-isophote light fraction
increases by $0.055$ in $V$ and $0.074$ in $g$. The supplement gives
per-box summaries and full paired aperture intervals.

\begin{figure*}
\centering
\includegraphics[width=\textwidth]{figures/validation_transfer.pdf}
\caption{External transfer test on 49 CAMELS Swift-EAGLE groups. Rows
show host-median overlap, absolute logarithmic luminosity error, and
absolute logarithmic projected half-light-radius error. The error rows
show $|\Delta\log_{10}L|$ and $|\Delta\log_{10}R_{50}|$, with
$\Delta\log_{10}L=\log_{10}(L_{\rm sel}/L_{\rm iso})$ and
$\Delta\log_{10}R_{50}=\log_{10}(R_{50,\rm sel}/R_{50,\rm iso})$.
Here, $L_{\rm sel}$ and $R_{50,\rm sel}$ refer to the selected group
or aperture, and $L_{\rm iso}$ and $R_{50,\rm iso}$ refer to the
isophotal region. Columns show $V$ and SDSS $g$. Each point is one host.
Box plots show the median and interquartile range. The x-axis categories
show $b_{\rm link}=0.10$, $b_{\rm link}=0.20$, and a 30 pkpc aperture.
The same simulation potential minimum centres all measurements.}
% Alt text: Six box-plot panels compare 49 CAMELS groups in V and g. Rows show overlap, absolute luminosity error, and absolute half-light-radius error for linking parameters 0.10 and 0.20 and a 30 pkpc aperture.
\label{fig:transfer}
\end{figure*}

\subsection{Test in separate CAMELS boxes}
\label{sec:camels-validation-results}

In the 130-host validation sample, fixed $b_{\mathrm{link}}=0.10$ raises mean
host-level IoU over $b_{\mathrm{link}}=0.20$ from 0.783 to 0.884 in $V$
and from 0.742 to 0.861 in $g$. The paired gains are 0.101
(95 per cent box-then-host interval 0.092--0.111) and 0.119
(0.108--0.130). Mean contamination falls by 0.140 and 0.153, while mean
completeness falls by 0.040 and 0.039. The 90th-percentile absolute
luminosity errors at $b_{\mathrm{link}}=0.10$ are 0.038 and 0.055 dex in
$V$ and $g$. The corresponding $R_{50}$ errors are 0.061 and 0.082 dex.
Large-failure fractions fall from 53.8 to 0 per cent in $V$ and from
76.2 to 2.3 per cent in $g$.

The low-redshift Darwin calibration ($b_{\mathrm{link}}=0.12$ in $V$ and
0.10 in $g$) also improves on 0.20 in both bands. The common fixed value
of 0.10 gives higher mean IoU than the Darwin value in $V$ (0.884 versus
0.864). Relative to a 30 pkpc aperture, fixed $b_{\mathrm{link}}=0.10$
raises mean IoU by 0.047 in $V$ (95 per cent interval 0.025--0.069) and
0.017 in $g$ ($-0.004$--$0.040$). The $g$ interval includes zero. The
fixed linking parameter lowers the 90th-percentile luminosity and
$R_{50}$ errors in both bands, but the aperture has lower contamination.
Median IoU remains below 0.90 for 67 per cent of $V$ hosts and 83 per
cent of $g$ hosts. The mean improvement therefore does not imply close
agreement for every host. The validation measures transfer across boxes within the Swift-EAGLE model,
not across simulation models.

\begin{table*}
\centering
\caption{Mean host-level metrics in the CV0--CV19 validation sample
($N=130$). Each host contributes the median across the three projections.
Large-failure fractions use the host-median definition in
Section~\ref{sec:camels-validation}. The fixed 30 pkpc aperture is an
unoptimized baseline. The $g$-band Darwin calibration is
$b_{\mathrm{link}}=0.10$ and is listed once.}
\label{tab:camels-validation}
\small
\begin{tabular}{llrrrrr}
\toprule
Band & Method & Setting & IoU & Completeness & Contamination & Large failures (\%) \\
\midrule
$V$ & FoF & 0.20 & 0.783 & 0.988 & 0.209 & 53.8 \\
$V$ & FoF & 0.10 & 0.884 & 0.948 & 0.070 & 0.0 \\
$V$ & FoF, Darwin calibration & 0.12 & 0.864 & 0.966 & 0.107 & 3.8 \\
$V$ & Projected aperture & 30 pkpc & 0.837 & 0.846 & 0.009 & 3.8 \\
$g$ & FoF & 0.20 & 0.742 & 0.989 & 0.251 & 76.2 \\
$g$ & FoF, Darwin calibration & 0.10 & 0.861 & 0.950 & 0.097 & 2.3 \\
$g$ & Projected aperture & 30 pkpc & 0.844 & 0.859 & 0.016 & 3.1 \\
\bottomrule
\end{tabular}
\end{table*}

\subsection{Transfer to TNG300-1}
\label{sec:tngresults}

The low-redshift Darwin parameters give median host-level IoU values of
0.892 in $V$ and 0.883 in $g$, compared with 0.833 and 0.783 at
$b_{\mathrm{link}}=0.20$. Median contamination is 0.056 and 0.053 at the
Darwin values, compared with 0.158 and 0.209 at 0.20. Paired host
bootstraps give median IoU gains of 0.067 in $V$ (95 per cent interval
0.055--0.073) and 0.104 in $g$ (0.093--0.109). The 90th-percentile
absolute luminosity errors are 0.031 and 0.042 dex, and the corresponding
projected $R_{50}$ errors are 0.057 and 0.065 dex. No host meets the
large-failure criterion at the Darwin parameters, compared with 22.5 per
cent in $V$ and 55 per cent in $g$ at 0.20. Median host-level IoU remains
below 0.90 for 65 per cent of $V$ hosts and 75 per cent of $g$ hosts.
The 30 pkpc aperture selects no light outside the isophote, while the
Darwin parameters give higher IoU and smaller luminosity and $R_{50}$
errors. The aperture comparison therefore retains a contamination
trade-off.

The TNG result depends strongly on the image centre. The proper distance
between catalogue positions \texttt{SubhaloCM} and \texttt{SubhaloPos}
has a median of 72.4 pkpc, a 90th percentile of 171.9 pkpc, and a maximum
of 396.7 pkpc. At fixed $b_{\mathrm{link}}=0.10$, the earlier
\texttt{SubhaloCM}-centred rendering has 90th-percentile absolute
luminosity errors of 3.36 dex in $V$ and 3.30 dex in $g$. The
\texttt{SubhaloPos}-centred rerender gives 0.059 and 0.042 dex. The
earlier stellar-density-peak-centred and present \texttt{SubhaloPos}
renderings agree within 0.002 in host-level median IoU for 90 per cent
of the hosts. Earlier TNG overlap measurements were inspected before we changed the
centre to \texttt{SubhaloPos}. The transfer test is therefore a post hoc,
fixed-parameter comparison in a second simulation model, based on one
snapshot. The single-output comparison does not establish general transfer across hydrodynamical
models.

\subsection{Image selection and inferred stellar growth}
\label{sec:growthbias}

\subsubsection{Growth rates and detection failures}

The central-subhalo stellar mass in the 36 common-central branches grows
by a median of $0.272$ dex from $z\simeq1$ to $z\simeq0.2$ (95 per cent
tree-bootstrap interval $0.218$--$0.368$ dex). When outputs are rendered
at the observer redshift of each simulation output, the $V$-band isophotal selection
exceeds central-subhalo stellar-mass growth by a median $0.176$ dex
($0.161$--$0.206$) across all 36 branches.

At $z\simeq1$, the central image pixel is brighter than the isophotal
threshold in all 108 $V$ projections and in 93 of 108 $g$ projections.
The other 15 $g$ projections use the region containing the nearest pixel
above the threshold. The fallback selection occurs in seven branches. We therefore treat
the 29 branches with central-pixel detections in every projection at
both endpoint outputs as the primary $g$-band comparison. Isophotal stellar-mass growth exceeds
central-subhalo growth by a median $0.196$ dex in the detected sample
(95 per cent interval $0.148$--$0.289$). Including the nearest-pixel
selections raises the all-branch median to $0.261$ dex
($0.182$--$0.303$). The fallback selections also change the mean bias from $0.210$ dex in
the detected sample to $0.526$ dex across all 36 branches because some
high-redshift fallback regions contain very few stars. The 16th--84th
percentile range of branch-level $g$-band growth errors is $0.069$--$0.691$
dex in all 36 branches and $0.054$--$0.323$ dex in the detection-qualified
sample. The
detection-qualified sample is selected by the image and may not represent
the full branch population. The $V$-band result is unchanged because all
36 branches have central-pixel detections at both endpoints.

The $g$-band growth comparison for every method uses both the all-branch
sample and the same 29-branch detection-qualified sample (Table~\ref{tab:growth-cohort}).
The calibrated FoF group at $b_{\mathrm{link}}=0.10$ has near-zero
median mass-growth bias in both samples. The $b_{\mathrm{link}}=0.20$
group also has near-zero median bias and the smaller mean absolute error.
The table therefore separates median bias from branch-level error and
shows the effect of applying the same detection cut to each method.

Rendering each physical output at the common observer redshift
$z_{\rm obs}=0.197$ changes the isophotal growth biases to $-0.070$ dex
in $V$ ($-0.087$ to $-0.062$) and $-0.079$ dex in $g$ ($-0.103$ to
$-0.073$). The paired difference, physical-observer minus fixed-observer
redshift, is $+0.266$ dex in $V$ ($+0.222$ to $+0.289$) and $+0.337$ dex
in $g$ ($+0.260$ to $+0.405$). The measured isophotal mass growth depends on image redshift as well as stellar
structure.

For $b_{\mathrm{link}}=0.10$, the median selected-to-central stellar-mass
ratios are $0.708$, $0.704$, and $0.708$ at $z\simeq1$, 0.5, and 0.2.
For $b_{\mathrm{link}}=0.20$, the ratios are $0.852$, $0.836$, and
$0.833$. The 16th--84th percentile ranges across branches are
$0.62$--$0.78$ for $b_{\mathrm{link}}=0.10$ and $0.77$--$0.88$ for
$0.20$. Both groups miss central-subhalo stellar mass, but similar
recovery fractions across epochs largely cancel in the growth difference.
The calibrated $b_{\mathrm{link}}=0.12$ in $V$ and $0.10$ in $g$ give
mean absolute mass-growth errors of $0.037$ and $0.044$ dex. The
$b_{\mathrm{link}}=0.20$ comparison gives $0.025$ dex in both bands.
Paired branch bootstraps give calibrated-minus-0.20 error differences of
$0.0115$ dex ($0.0040$--$0.0197$) in $V$ and $0.0184$ dex
($0.0094$--$0.0280$) in $g$. The projected 30 pkpc aperture has a
median mass-growth bias of $-0.127$ dex; the Darwin-calibrated spherical
apertures give $-0.067$ dex in $V$ and $-0.088$ dex in $g$.

The linking length is defined in comoving coordinates (equation~\ref{eq:link}).
At fixed $b_{\mathrm{link}}$ and particle mass, the proper linking length at $z=0.997$
is $0.60$ of the linking length at $z=0.197$. Over the same interval, the central
half-mass radius grows by $0.382$ dex, or a factor of $2.41$. Thus the nearly
constant selected-to-central mass ratios do not correspond to a fixed proper
aperture. The recovered mass fractions alone do not explain why the recovered fractions
remain similar.

The isophotal projected half-light-radius growth bias relative to stars
assigned to the central subhalo is a median $0.430$ dex in $V$ and
$0.296$ dex in $g$ for the full sample; the $g$-band median is $0.334$
dex in the detection-qualified sample. For stars bound to the
central subhalo, half-light-radius growth is $0.020$ dex in $V$ and
$-0.003$ dex in $g$ when each output is rendered at the observer redshift of the simulation output; both 95 per cent intervals include zero. At fixed
observer redshift, the corresponding medians are $0.257$ and $0.237$
dex. The central-subhalo projected half-mass radius grows by $0.382$
dex ($0.306$--$0.450$) under both observer conditions. The same
central-subhalo membership, proper positions, and full particle fluxes
define the reference radius in each pair of renderings. Surface-brightness
dimming, PSF convolution, and pixel scale do not enter the reference
$R_{50}$. A common luminosity-distance factor also cancels from a
half-light radius. The change between observer conditions therefore
arises from spectral reweighting: the observed filters sample different
rest-frame wavelengths, changing the relative fluxes of stars with
different spectra. Stellar-population evolution occurs in both
comparisons. Image selection adds separate changes to the isophotal
$R_{50}$. The central-bound half-mass radius grows by $0.382$ dex,
whereas the calibrated FoF groups have half-mass-radius growth biases
consistent with zero. Half-light-radius growth
depends on luminosity weighting and image condition; half-light radius
alone does not measure structural growth. The 36 branches come from one
simulation volume and do not establish universal corrections.

\subsubsection{Changes in central-subhalo mass}

Tracking exact stellar-particle IDs separates changes in central-subhalo
membership from changes in stellar mass. From $z\simeq1$ to $z\simeq0.2$, the median
central stellar-mass increase is $0.869$ times the starting central mass.
Pre-existing stars newly assigned to the central contribute $0.789$ times
the starting mass; stars formed after the first output and assigned to the
central at the last output contribute $0.096$. The current masses of
persistent central-member particles change by $-0.050$ times the starting
mass, and particles leaving central membership contribute $-0.002$.
The component contributions are separate branch medians and need not sum to the median net
increase. The mass balance closes branch by branch, with a maximum fractional
residual of $9.7\times10^{-16}$. Of the pre-existing stars newly assigned
to the central, a median 88 per cent carried the resolved-satellite mask at
$z\simeq1$. The mask identifies an earlier catalogue class. The satellite mask does not
identify a donor subhalo or prove a merger origin. The unflagged remainder
is 11.9 per cent of the newly acquired stellar mass in the branch-level medians: 2.1 per cent was
in the earlier parent FoF but was not assigned to a resolved satellite, and
8.2 per cent was absent from the earlier parent-FoF catalogue. The
component medians need not sum to the median remainder; the earlier location
of the absent stars is unknown.

\subsubsection{Stellar-population weighting of half-light radius}

We also repeat the central-subhalo half-light calculation after replacing
each particle's age and total metal mass fraction by the central population's
current-mass-weighted geometric means at each output. Each geometric mean
is the exponential of the current-mass-weighted mean logarithm. The uniform-population control
removes particle-to-particle differences in age and metallicity weights,
including spatial gradients. The control does not isolate radial gradients
alone. With spatially uniform population weights, half-light-radius growth
is $0.386$ dex in both bands. The uniform-population half-light radius
weights particles by formed stellar mass, whereas the half-mass radius uses
current particle mass and grows by $0.382$ dex. The paired change in
half-light-radius growth (full particle-level age and metallicity weights
minus uniform weights) is $-0.255$ dex in $V$ (95 per cent interval
$-0.440$ to $-0.175$) and $-0.324$ dex in $g$ ($-0.441$ to $-0.225$)
for images rendered at the observer redshift of each simulation output. At $z\simeq1$,
stars outside the projected
central half-mass radius are younger by $0.30$ Gyr and lower in log total
metal mass fraction by $0.51$ dex than stars inside the half-mass radius. The outer-minus-inner
log mass-to-light ratio differences are $-0.556$ dex in $V$ and $-0.775$ dex
in $g$ at the physical observer redshift, compared with $-0.245$ and $-0.301$
dex when both outputs are rendered at $z_{\rm obs}=0.197$. At $z\simeq0.2$,
the differences are $-0.168$ dex in $V$ and $-0.205$ dex in $g$ under either
observer condition. The stronger high-redshift contrast occurs because the
observed bands sample rest-frame near-UV light. The paired full-population
half-light-radius growth is lower by $0.183$ dex in $V$ and $0.199$ dex in
$g$ at the physical observer redshifts than at fixed $z_{\rm obs}=0.197$.
For the full-minus-uniform growth change, the age and metallicity Shapley
contributions at the physical observer redshifts are $-0.162$ and $-0.123$
dex in $V$, and $-0.175$ and $-0.143$ dex in $g$. The component medians need not add
to the median combined change. The mock images omit dust, and the dust-free
results do not determine the size of the effect in observed galaxies.

\begin{table*}
\centering
\caption{Physical-redshift $g$-band stellar-mass growth relative to
stars bound to the central subhalo. Each branch contributes the median
over three views. The detection-qualified sample contains the 29 branches
with central-pixel detections in every view at both endpoints. Bias is
the median branch difference in growth, and MAE is the mean of the
absolute branch differences. The Supplement lists both photometric bands
and all 36 branches.}
\label{tab:growth-cohort}
\small
\begin{tabular}{lrrrr}
\toprule
Selection & Bias, all 36 & MAE, all 36 & Bias, detected 29 & MAE, detected 29 \\
\midrule
Isophotal region & $+0.261$ & 0.531 & $+0.196$ & 0.217 \\
FoF, $b_{\mathrm{link}}=0.10$ & $-0.004$ & 0.044 & $-0.005$ & 0.047 \\
FoF, $b_{\mathrm{link}}=0.20$ & $-0.003$ & 0.025 & $-0.004$ & 0.026 \\
Projected aperture, 30 pkpc & $-0.127$ & 0.125 & $-0.125$ & 0.123 \\
Spherical aperture, 60 pkpc & $-0.088$ & 0.098 & $-0.087$ & 0.098 \\
\bottomrule
\end{tabular}
\end{table*}

\begin{table*}
\centering
\caption{The preferred linking parameter depends on the measurement
objective. The full-image calibration maximizes light-weighted overlap.
Boundary-band overlap and central-subhalo mass-growth error are evaluated
separately. The boundary result depends on band width.}
\label{tab:objectives}
\small
\begin{tabular}{p{3.1cm}p{4.0cm}p{7.0cm}}
\toprule
Objective & Sample and comparison & Result ($V/g$) \\
\midrule
Full-image stellar-light agreement &
130-host CV0--CV19 validation sample; fixed $b_{\mathrm{link}}=0.10$ versus 0.20 &
Mean host IoU increases by $0.101/0.119$; completeness decreases by
$0.040/0.039$. The Darwin full-sample medians are
$b_{\mathrm{link}}=0.08/0.10$ on the original nine-point grid. \\
Near-isophote overlap at $\Delta=0.5$ arcsec &
75 retained Darwin entries; pixel-centre distance threshold; original nine-point positional calibration versus 0.20 &
Median boundary-band IoU is $0.516/0.528$ at the selected values and
$0.575/0.592$ at 0.20. The ranking reverses at a $2$-arcsec threshold. \\
Central-subhalo stellar-mass growth &
36 TNG300-1 branches; calibrated values ($0.12/0.10$) versus 0.20 &
Mean absolute growth error is $0.037/0.044$ dex at the calibrated values
and $0.025/0.025$ dex at 0.20. Paired calibrated-minus-0.20 differences
are $0.0115$ dex ($0.0040$--$0.0197$) in $V$ and $0.0184$ dex
($0.0094$--$0.0280$) in $g$. \\
\bottomrule
\end{tabular}
\end{table*}

\begin{figure*}
\centering
\includegraphics[width=0.92\textwidth]{figures/tng300_tree_growth_bias_r1r4d.pdf}
\caption{Bias in stellar-mass growth (top) and projected half-light-radius
growth (bottom), defined as the growth measured by each selection minus
the growth of stars bound to the central subhalo. Filled symbols use the
physical observer redshift of each output; open symbols render all outputs
at $z_{\rm obs}=0.197$. Blue and orange symbols show $V$ and $g$,
respectively. Circles show medians across all 36 branches. The orange
diamond at the isophotal selection shows the $g$-band result for the 29
branches with central-pixel detections in all three projections at both
endpoints. Error bars show 95 per cent bootstrap intervals with the branch
as the sampling unit. The vertical axes give $\Delta\log_{10}M_\star$ and
$\Delta\log_{10}R_{50}$, the selected-minus-central-subhalo growth from
$z\simeq1$ to 0.2, in dex. Positive values mean that the selection gives
greater growth. The FoF comparisons include $b_{\mathrm{link}}=0.10$,
0.12, and 0.20; 0.12 in $V$ and 0.10 in $g$ are the calibrated values.
The isophotal circles use all 36 branches and include
regions selected around the nearest pixel above the threshold;
the diamond shows the detection-qualified primary result. Table~\ref{tab:growth-cohort}
compares all methods in the full and detection-qualified $g$-band samples.
The spherical-aperture radii are
80 pkpc in $V$ and 60 pkpc in $g$.}
% Alt text: Two vertically joined panels compare stellar-mass and half-light-radius growth bias for six selections. Circles show all 36 TNG branches; an orange diamond marks the g-band isophotal result for 29 branches with central-pixel detections in every view at both endpoints. Filled circles use physical observer redshifts and open circles use a fixed observer redshift.
\label{fig:tng-growth-bias}
\end{figure*}

\subsection{Artificial redshift and image construction}
\label{sec:artificialredshiftresults}

For observer-frame images at fixed $b_{\mathrm{link}}=0.10$, the mean overlap across
all 225 views at $z_{\rm obs}=0.190$, $0.227$, $0.500$, $1.002$, and
$1.990$ is $0.902$, $0.907$, $0.897$, $0.816$, and $0.671$ in $V$, and
$0.906$, $0.902$, $0.895$, $0.817$, and $0.719$ in $g$. Views without a
detected isophotal region contribute zero overlap. The 75 halo-output
entries contribute three views each. Entries at different outputs can be
repeated systems. The means describe the sample and have no confidence
intervals based on independent haloes. The sample includes outputs at
$z\leq1.990$, omitting the single $z=3.002$ entry and including the
$z=0.227$ halo excluded from calibration because the maximum overlap
occurs at a grid endpoint.

At $z_{\rm obs}=1.990$, an isophotal region is detected in 182/225
$V$ views and 188/225 $g$ views. Among detected views, mean completeness
is $0.996$ in $V$ and $0.994$ in $g$, and mean overlap is $0.829$ and
$0.860$. Assigning zero overlap to views without a detected region gives
means of $0.671$ and $0.719$ across all views. Removing the $(1+z)^4$
cosmological-dimming factor while retaining the redshifted spectrum gives
all-view means of $0.881$ in $V$ and $0.877$ in $g$, increases of $0.211$
and $0.158$. Replacing the redshifted spectrum with the $z=0$ spectrum
while retaining the Darwin luminosity distance gives $0.795$ and $0.818$.
In the extreme no-distance control, the isophote fills the full 360-arcsec
stamp in 2 of 225 views in each band. The mean overlap depends on the
field of view and is not used as a physical comparison. The pixel scale and PSF are fixed
in angle, so physical resolution also changes with rendering redshift. The
comparisons measure combined image-construction effects, not physical
halo evolution. Figure~\ref{fig:artificial} shows the mean overlap as a function of observer redshift.
Table~\ref{tab:artificial-detection} lists the detection counts and mean overlaps at $z_{\rm obs}=1.990$.
\begin{figure}
\centering
\includegraphics[width=\columnwidth]{figures/artificial_redshift.pdf}
\caption{Mean light-weighted overlap across all 225 projections as a function
of observer redshift at fixed $b_{\mathrm{link}}=0.10$. Projections without
a detected isophotal region receive zero overlap. Stellar positions and
populations remain fixed. The angular pixel scale and PSF remain fixed,
so the physical resolution changes with redshift.}
\label{fig:artificial}
\end{figure}

\begin{table}
\centering
\caption{Isophotal detections and mean overlap at $z_{\rm obs}=1.990$.
The detected-view mean excludes non-detections. The all-view mean assigns
zero overlap to non-detections. Each band contains 225 projections.}
\label{tab:artificial-detection}
\begin{tabular}{lrrr}
\hline
Band & Detections & $\langle\mathrm{IoU}\rangle_{\rm det}$ & $\langle\mathrm{IoU}\rangle_{\rm all}$ \\
\hline
$V$ & 182/225 & 0.829 & 0.671 \\
$g$ & 188/225 & 0.860 & 0.719 \\
\hline
\end{tabular}
\end{table}

We repeat the fixed-$b_{\mathrm{link}}$ comparison for 53 entries from outputs
$z=1.990$, $1.002$, and $0.500$. Twenty major-descendant links group the entries
into 33 components. We reduce the three views at each output before taking the
median across repeated outputs in each component. At $z_{\rm obs}=0.190$ and
$0.227$, the paired mean IoU change between $b_{\mathrm{link}}=0.10$ and
$0.20$ is consistent with zero in both bands. At $z_{\rm obs}=0.500$,
$1.002$, and $1.990$, $b_{\mathrm{link}}=0.10$ raises mean IoU by
$0.039$, $0.039$, and $0.021$ in $V$, and by $0.031$, $0.029$, and $0.022$
in $g$. All six 95 per cent group-bootstrap intervals exclude zero
(supplementary table). At $z_{\rm obs}=1.990$, 24 per cent of components in
$V$ and 21 per cent in $g$ remain below IoU $0.90$. The 33 components use
major exact-ID links and do not form a complete merger tree. Unidentified
links can leave repeated systems in separate components.

A separate fixed-structure test renders 49 CAMELS Swift-EAGLE hosts at four
observer redshifts and five surface-brightness-limit combinations. At the
nominal limits and fixed $b_{\mathrm{link}}=0.10$, the mean host overlap falls from
$0.894/0.872$ at $z_{\rm obs}=0.190$ to $0.690/0.653$ at $0.500$,
$0.263/0.300$ at $1.002$, and $0.041/0.116$ at $1.990$ ($V/g$).
A contour cannot be measured in at least one of three views for 29/32 hosts
at $z_{\rm obs}=1.002$ and 46/44 hosts at $1.990$. A fixed 30 pkpc aperture
has higher mean overlap than $b_{\mathrm{link}}=0.10$ at $z_{\rm obs}=0.500$ and
$1.002$. A 0.5-mag deeper limit improves overlap but still leaves failed views
in 43 $V$-band and 41 $g$-band hosts at $z_{\rm obs}=1.990$. The rendering
holds stellar structure and group membership fixed while changing redshifted
spectra, cosmological dimming, and angular resolution together. The test
limits image-measurement transfer and does not represent physical galaxy
evolution; the supplementary material gives the threshold results and paired
comparisons.

\subsection{Agreement near the brightness boundary}
\label{sec:residual}

Stars near the galaxy centre supply much of the light used to calculate
overlap. The ratio can therefore remain close to unity even when group
members and stars inside the isophote differ near the boundary.
For a comparison near the boundary, we assign each image pixel a signed
Euclidean distance to the centre of the nearest pixel of the opposite
isophotal class. The distance is negative inside the contour and positive
outside. At threshold $\Delta$, the boundary band contains pixels with
$|d|\leq\Delta$; the distance is between pixel centres, not from a
continuous contour. We evaluate five thresholds from $0.2$ to $2.0$ arcsec.
The $0.5$-arcsec threshold is comparable to the $1$-arcsec PSF full width
at half maximum (FWHM), but does not define an exact one-FWHM-wide band.
Equation~\ref{eq:iou} then uses only stars projected into these pixels for
all three luminosity sums. Group members farther from the boundary
contribute to neither the numerator nor the denominator. The angular
threshold spans different physical distances at different redshifts. The
original one-pixel ring, including pixels touching at a corner, corresponds
to a $0.3$-arcsec threshold on this image grid. At $\Delta=0.3$ arcsec,
median overlap is $0.461$ in $V$ and $0.476$ in $g$ at
the selected parameters, compared with $0.541$ and $0.548$ at
$b_{\mathrm{link}}=0.20$. At $\Delta=0.5$ arcsec, the corresponding
medians are $0.516/0.528$ and
$0.575/0.592$ ($V/g$). At a $2$-arcsec distance threshold, the ranking reverses:
the selected parameters give $0.842/0.862$, compared with $0.808/0.851$
at $b_{\mathrm{link}}=0.20$. Thus boundary agreement depends on the
spatial scale of the comparison, and the one-pixel result does not by
itself measure general boundary accuracy.

Archived-normalization fractions of group light outside and non-member light inside the contour are tabulated in Supplementary Table S1. The archived fractions do not enter the corrected-distance results.

\subsection{Where the match fails}
\label{sec:fail}

Nine retained haloes have median overlap below $0.90$ in both bands.
The weakest cases arise when the isophotal region covers only a few
pixels or does not contain the stellar-density seed. In one $z=1.002$
$g$-band view, the archived normalization selects one pixel of
$0.04\,\mathrm{arcsec}^{2}$, giving IoU $0.0042$. With the consistent
Darwin distance, the selected area is $17.16\,\mathrm{arcsec}^{2}$ and
the maximum IoU is $0.9073$ at $b_{\mathrm{link}}=0.10$; the change
occurs because a different bright region is selected. At $z=0.227$, one
halo has an endpoint maximum in two views and is excluded from the
calibration. Its image centre lies more than $30\,\mathrm{arcsec}$ from
the density seed. Selecting the seed-containing region raises the
overlap, but the primary analysis retains the centre-based image rule.
Across the 456 band-view combinations, 450 selected regions already
contain the density seed; using the seed-containing region changes no
selected parameter or median overlap for the 75 retained entries.
Supplementary Fig. S2 preserves the archived-normalization maps. The
examples show sensitivity to the image-region rule and photometric
normalization; they do not establish a physical galaxy boundary.

An archived control without cosmological dimming and with unredshifted spectra continues to favour larger groups through the tested $b_{\mathrm{link}}=0.50$ limit. The control changes both the sampled wavelength and the surface-brightness normalization, so it is not used to infer a redshift trend or a primary calibration; the full method and curves appear in Supplementary Fig. S3.

\section{Discussion}
\label{sec:use}
\subsection{Objective-dependent interpretation}

Photometric and dynamical definitions
of diffuse cluster light select different stars
\citep{tang2018,montenegro2025}. Our analysis quantifies the consequences
for three measurements, total light inside an isophote, light near the
boundary, and stellar-mass growth relative to central-subhalo members.
The Darwin-calibrated grouping improves
full-image overlap in additional boxes from the same simulation model. The additional boxes were not used to set the linking parameter.
The same calibration does
not optimize boundary overlap and central-subhalo growth (Table~\ref{tab:objectives}).
The comparison identifies the selection giving the smallest error for
each measured quantity in the specified synthetic images. The comparison
does not establish a new physical mechanism or a universal definition of
intracluster light.

\subsection{Image construction and particle sampling}
\label{sec:sensitivity}

We keep $b_{\mathrm{link}}=0.10$ and $0.20$ fixed while varying one image choice for
all 49 CAMELS hosts. The checks change the brightness limits by
$\pm0.5$ mag, double the pixel size, remove the PSF wing, add a flat
sky residual, or add Gaussian pixel noise. No linking parameter is
retuned. We take the median over three views within each noise
realisation and average the three realisations within each host.
Noiseless conditions use one realisation. The supplementary notes
give the noise amplitudes, failure accounting, and full tables.

In every tested condition and both bands, $b_{\mathrm{link}}=0.10$ gives a higher mean
overlap than $b_{\mathrm{link}}=0.20$, with positive 95 per cent box-then-host paired
bootstrap intervals. When the pixel noise rms equals the flux at the
brightness limit, the mean overlaps at $b_{\mathrm{link}}=0.10$ are $0.877$ in $V$
and $0.857$ in $g$. No central isophote fails detection in the tested
conditions. Changing the brightness limit shifts the mean signed
luminosity bias from $+0.006/+0.017$ dex in the baseline images to
$+0.039/+0.055$ dex for brighter limits and $-0.023/-0.015$ dex for
fainter limits ($V/g$). In $V$, fainter limits raise mean overlap from
$0.894$ to $0.901$ but increase the 90th-percentile luminosity error
from $0.033$ to $0.056$ dex. Higher overlap therefore need not reduce
the largest luminosity errors.

The image perturbations change the isophotal boundary. Luminosity and
radius measurements retain noiseless particle fluxes, so the checks
do not test noisy aperture photometry. Dust, correlated sky errors,
and image-based centre finding remain untested.
Separate Darwin particle-thinning checks retain the archived flux
normalization. With a fixed image, retaining one half, one quarter,
or one eighth of the linking particles raises the median selected $b_{\mathrm{link}}$
by $0.02$, $0.05$, or $0.07$. The one-eighth result retains only
30 of 75 haloes per band after excluding endpoint maxima. The
supplementary notes preserve the archived image and sampling checks.
Particle thinning does not establish simulation convergence.

\subsection{Limits of the calibration}

The centre pixel is fainter than the limit in 39 $V$ projections and 36
$g$ projections. The nearest above-threshold region is selected in these
views; 31 $V$ and 28 $g$ views have positional overlap of at least $0.90$
at the parameters selected from the full images. Requiring centre-pixel detections
in all three projections leaves 55 $V$ and 53 $g$ entries, with unchanged
median selected parameters. Thus the fallback rule affects individual
selections but does not shift the median calibration in the detection-selection test.
The Darwin calibration has not been tested with interstellar dust,
additional viewing directions, or fixed $30$ or $100\,\mathrm{kpc}$ apertures. The projected 30 pkpc aperture is
included in the CAMELS transfer test and the TNG300-1 growth comparison.
The projected-aperture comparison uses a fixed 30 pkpc radius and does
not establish a general FoF advantage over a separately calibrated
projected aperture. The TNG growth test includes Darwin-calibrated
three-dimensional sphere radii, but the different geometry does not
replace a matched projected-aperture test.
The CV0--CV19 validation sample covers additional boxes from the Swift-EAGLE model,
while TNG300-1 supplies a post hoc transfer comparison and a separate
merger-tree growth analysis from one simulation volume. The 40-host
centre-sensitive transfer test uses one output; the growth analysis follows
36 branches across three outputs. Earlier
\texttt{SubhaloCM}-centred and current \texttt{SubhaloPos}-centred
renderings give substantially different luminosity errors. The difference
shows sensitivity to image centring and limits claims of model transfer.
The artificial-redshift tests keep stellar structure fixed and combine
changes in image construction; the tests do not measure physical redshift
evolution.


\section{Conclusions}
\label{sec:sum}

We calibrate three-dimensional stellar FoF groups against the stellar
light inside specified surface-brightness contours. On the original
nine-point grid, the median selected $b_{\mathrm{link}}$ across 75 retained
Darwin halo-output entries is 0.08 in $V$ and 0.10 in $g$, with 16th--84th
percentile ranges of 0.08--0.12 in both bands. The parameter maximizes
full-image light-weighted overlap for the tested image construction and
sample; the value does not define a universal stellar boundary.
On the corrected-distance 15-point grid, the median individual optimum is
$0.09$ in both bands for positional and PSF-aperture weighting. The
median-over-entry positional curves peak at $0.09$ in both bands, and
the PSF-aperture curves at $0.09$ in $V$ and $0.10$ in $g$. A fixed
$b_{\mathrm{link}}=0.10$ is therefore a practical choice for the stated
integrated-light comparison across the tested mock-image samples, rather than
a universal optimum.

A test that omits one halo at a time at $z=0.1904$ selects 0.12 in $V$
and 0.10 in $g$ in every case; all 12 omitted haloes have higher overlap
than at 0.20.
In the 130-host CV0--CV19 validation sample, $b_{\mathrm{link}}=0.10$
raises mean host IoU over 0.20 by 0.101 in $V$ and 0.119 in $g$.
Contamination falls by 0.140 and 0.153, while completeness falls by
0.040 and 0.039. The validation sample supports transfer across boxes within the
Swift-EAGLE model. The fixed 30 pkpc aperture is an unoptimized baseline,
so the comparison does not establish a general FoF advantage over an
aperture radius calibrated on a separate sample.

The preferred selection depends on the measurement. At a pixel-centre
distance threshold of $0.5$ arcsec around the contour, the full-image-
calibrated values give lower overlap than
$b_{\mathrm{link}}=0.20$ (median IoU $0.516/0.528$ versus $0.575/0.592$
in $V/g$); at a $2$-arcsec threshold, the ranking reverses. In TNG300-1, the
calibrated linking values give near-zero median
bias in central-subhalo mass growth, but $b_{\mathrm{link}}=0.20$ gives
smaller mean absolute errors: 0.025 dex in both bands, compared with
0.037 dex in $V$ and 0.044 dex in $g$ at the calibrated values. The
paired differences are positive in both bands. The calibration therefore
serves the stated integrated-light objective; neither boundary overlap
nor central-subhalo growth is optimized by that choice.

In 36 TNG300-1 branches, isophotal mass growth exceeds central-subhalo
growth by a median 0.176 dex in $V$. For $g$, the median is 0.196 dex
among 29 branches with central-pixel detections in every view at both
endpoints. Fifteen of 108 high-redshift $g$ projections select the
region around the nearest pixel above the threshold, across seven
branches; including those projections raises the all-branch median to
0.261 dex. The detection-qualified
sample is selected by the image and may carry selection effects.
FoF groups at both 0.10 and 0.20 miss central-subhalo mass at every
epoch, while relatively stable recovery fractions produce small growth
biases. Particle-ID tracking assigns a median $0.789$ times the initial
central mass to pre-existing stars newly classified as central members;
88 per cent of this mass carried the resolved-satellite mask at snapshot 50
($z=0.997$), the earlier endpoint of the snapshot-50-to-84 comparison.
The catalogue mask cannot identify a donor or establish a
physical merger origin. The 36 branches come from one volume and do not
provide a universal correction.

For central-subhalo stars at fixed snapshot, the half-light reference
uses unchanged membership and proper positions in each observer
condition. Its median growth is 0.020 dex in $V$ and $-0.003$ dex in
$g$ at physical observer redshifts, compared with 0.257 and 0.237 dex
at fixed $z_{\rm obs}=0.197$. The difference arises from spectral
reweighting of stars with different spectra, not from changes in
membership or position. The population control changes the luminosity weights
of individual stars by age and metallicity, including spatial gradients.
it does not separate a radial-gradient effect from other particle-to-particle
population differences. The projected half-mass radius grows by
0.382 dex. In controlled dust-free SSP renderings, making stellar ages
and metallicities uniform raises half-light-radius growth by a paired
median 0.255 dex in $V$ and 0.324 dex in $g$ at physical observer
redshifts. The outer stars are younger and more metal-poor than the inner
stars, and have lower mass-to-light ratios, especially at $z\simeq1$.
Half-light-radius growth therefore depends on observer condition and
spatially varying stellar populations; it does not by itself measure
structural growth.

The CAMELS validation sample uses one simulation model. The TNG comparison uses
one finite-resolution volume, with no resolution-convergence test.
Artificial-redshift images hold stellar structure fixed and change
image construction, so the tests do not measure physical evolution.
The results calibrate a projected image selection; they do not identify
gravitationally unbound stars or measure the intracluster-light fraction.
\FloatBarrier
\section*{Acknowledgements}

J. Lee acknowledges support from the National Research Foundation of Korea
(NRF) grant funded by the Korean government (MSIT; RS2022-NR068800).
Juhan Kim acknowledges support from a KIAS Individual Grant (KG039603) at
the Korea Institute for Advanced Study. Large data transfers were supported
by KREONET, managed and operated by the Korea Institute of Science and
Technology Information. The authors thank the Center for Advanced Computation
at the Korea Institute for Advanced Study for computational resources.
Juhan Kim provided the MPI+OpenMP halo-finding and galaxy-finding codes used
in this work. OpenAI Codex and Anthropic Claude assisted with code development,
additional analyses, figure preparation, and manuscript editing. The authors reviewed and edited the AI-generated text and are responsible
for the scientific claims and the final manuscript.

\section*{Competing interests}
The authors declare no competing interests.

\section*{Data availability}

The author-review release PGalF-MNRAS-editorial-20261107-r5 contains the
manuscript source, supplementary tables, analysis scripts, figure inputs,
and checksums. Co-authors may request the local package from the
corresponding author. The manuscript review copy is maintained on Overleaf. The analysis source
code is available at \url{https://github.com/kjhan0606/GalaxyFinder}.
The complete local data package has no public DOI.
REPRODUCIBILITY.md gives the verification and build commands. The existing
Overleaf copy omits the binary TNG300-1 cutouts and compiled linking library;
the local review package includes them.

The original Darwin particle catalogues are not included. Reconstructing
the stellar memberships and rerunning the primary image measurements
requires access to those simulation inputs. Requests for access should be
directed to the corresponding author. The package includes the archived
stellar-population grids and the derived Darwin overlap measurements. The
local review package includes the TNG300-1 cutouts, SubLink branches, growth
measurements, and source dependencies for the reported sample. The CV0--CV19 CAMELS validation sample
measurements and scripts, TNG300-1 image-transfer metrics, and supplementary
sensitivity tables are included in the corresponding analysis directories. Full
reconstruction may require access credentials for TNG300-1 data and
approval to use the Darwin catalogues.

\begin{thebibliography}{99}
\bibitem[Brough et al.(2024)]{brough2024}
Brough, S., Ahad, S.~L., Bah\'e, Y.~M., et al. 2024, MNRAS, 528, 771
\bibitem[Burke et al.(2015)]{burke2015}
Burke, C., Hilton, M., \& Collins, C. 2015, MNRAS, 449, 2353
\bibitem[Conroy et al.(2009)]{conroy2009}
Conroy, C., Gunn, J.~E., \& White, M. 2009, ApJ, 699, 486
\bibitem[Contini(2021)]{contini2021}
Contini, E. 2021, Galaxies, 9, 60
\bibitem[Cui et al.(2014)]{cui2014}
Cui, W., Murante, G., Monaco, P., et al. 2014, MNRAS, 437, 816
\bibitem[Davis et al.(1985)]{davis1985}
Davis, M., Efstathiou, G., Frenk, C.~S., \& White, S.~D.~M. 1985, ApJ, 292, 371
\bibitem[Feldmeier et al.(2004)]{feldmeier2004}
Feldmeier, J.~J., Mihos, J.~C., Morrison, H.~L., Harding, P., Kaib, N., \& Dubinski, J. 2004, ApJ, 609, 617
\bibitem[Gonzalez et al.(2005)]{gonzalez2005}
Gonzalez, A.~H., Zabludoff, A.~I., \& Zaritsky, D. 2005, ApJ, 618, 195
\bibitem[Gonzalez et al.(2007)]{gonzalez2007}
Gonzalez, A.~H., Zaritsky, D., \& Zabludoff, A.~I. 2007, ApJ, 666, 147
\bibitem[Kim et al.(2023)]{kim2023}
Kim, J., Lee, J., Laigle, C., et al. 2023, ApJ, 951, 137
\bibitem[Kluge et al.(2020)]{kluge2020}
Kluge, M., Neureiter, B., Riffeser, A., et al. 2020, ApJS, 247, 43
\bibitem[Marinacci et al.(2018)]{marinacci2018tng}
Marinacci, F., et al. 2018, MNRAS, 480, 5113
\bibitem[Mihos et al.(2005)]{mihos2005}
Mihos, J.~C., Harding, P., Feldmeier, J., \& Morrison, H. 2005, ApJ, 631, L41
\bibitem[Mihos et al.(2017)]{mihos2017}
Mihos, J.~C., Harding, P., Feldmeier, J.~J., et al. 2017, ApJ, 834, 16
\bibitem[Montenegro-Taborda et al.(2023)]{montenegro2023}
Montenegro-Taborda, D., Rodriguez-Gomez, V., Pillepich, A., et al. 2023, MNRAS, 521, 800
\bibitem[Montenegro-Taborda et al.(2025)]{montenegro2025}
Montenegro-Taborda, D., Rodriguez-Gomez, V., Avila-Reese, V., et al. 2025, MNRAS, 541, 2162
\bibitem[Montes(2022)]{montes2022}
Montes, M. 2022, Nat.\ Astron., 6, 308
\bibitem[More et al.(2011)]{more2011}
More, S., Kravtsov, A.~V., Dalal, N., \& Gottl\"ober, S. 2011, ApJS, 195, 4
\bibitem[Naiman et al.(2018)]{naiman2018tng}
Naiman, J.~P., et al. 2018, MNRAS, 477, 1206
\bibitem[Nelson et al.(2018)]{nelson2018tng}
Nelson, D., et al. 2018, MNRAS, 475, 624
\bibitem[Nelson et al.(2019)]{nelson2019tng}
Nelson, D., et al. 2019, Comput.\ Astrophys.\ Cosmol., 6, 2
\bibitem[Pillepich et al.(2018)]{pillepich2018tng}
Pillepich, A., et al. 2018, MNRAS, 475, 648
\bibitem[Rodriguez-Gomez et al.(2015)]{rodriguezgomez2015}
Rodriguez-Gomez, V., Genel, S., Vogelsberger, M., et al. 2015, MNRAS, 449, 49
\bibitem[Rudick et al.(2011)]{rudick2011}
Rudick, C.~S., Mihos, J.~C., \& McBride, C.~K. 2011, ApJ, 732, 48
\bibitem[Springel et al.(2018)]{springel2018tng}
Springel, V., et al. 2018, MNRAS, 475, 676
\bibitem[Tang et al.(2018)]{tang2018}
Tang, L., Lin, W., Cui, W., et al. 2018, ApJ, 859, 85
\bibitem[Teyssier(2002)]{teyssier2002}
Teyssier, R. 2002, A\&A, 385, 337
\bibitem[Tolman(1930)]{tolman1930}
Tolman, R.~C. 1930, Proc.\ Natl.\ Acad.\ Sci., 16, 511
\bibitem[Villaescusa-Navarro et al.(2021)]{villaescusa2021}
Villaescusa-Navarro, F., Angl\'es-Alc\'azar, D., Genel, S., et al. 2021, ApJ, 915, 71
\bibitem[Zhang et al.(2019)]{zhang2019}
Zhang, Y., Yanny, B., Palmese, A., et al. 2019, ApJ, 874, 165
\bibitem[Zibetti et al.(2005)]{zibetti2005}
Zibetti, S., White, S.~D.~M., Schneider, D.~P., \& Brinkmann, J. 2005, MNRAS, 358, 949
\bibitem[Zwicky(1951)]{zwicky1951}
Zwicky, F. 1951, PASP, 63, 61
\end{thebibliography}

\bsp
\label{lastpage}
\end{document}
